\begin{problem}[33届物理竞赛复赛]{58}
    木星是太阳系内质量最大的行星（其质量约为地球的 318 倍）。假设地球与木星均沿圆轨道绕太阳转动，两条轨道在同一平面内。将太阳、地球和木星都视为质点，忽略太阳系内其它星体的引力；且地球和木星之间的引力在有太阳时可忽略。已知太阳和木星质量分别为 $m_{s}$ 和 $m_{j}$ ，引力常量为 $G$ 。地球和木星绕太阳运行的轨道半径分别是 $r_{e}$ 和 $r_{j}$ 。假设在某个时刻，地球与太阳的连线和木星与太阳的连线之间的夹角为 $\theta$ 。这时若太阳质量突然变为零，求

    \begin{subproblem}
        此时地球相对木星的速度大小 $v_{ej}$ 和地球不被木星引力俘获所需要的最小速率 $v_{0}$ 。
    \end{subproblem}
    \begin{subproblem}
        试讨论此后地球是否会围绕木星转动，可利用（1）中结果和数据 $m_{\mathrm{s}} \approx 2.0 \times 10^{30} \mathrm{~kg}$ 、 $m_{\mathrm{j}} \approx 1.9 \times 10^{27} \mathrm{~kg}$ 、木星公转周期 $T_{\mathrm{j}} \approx 12 \mathrm{~y}$ 。
    \end{subproblem}
\end{problem}